\MulaiModul{11}{Debugging}

\subsection{Tujuan}

\begin{tujuan}
Setelah menyelesaikan modul ini, mahasiswa mampu:
\begin{enumerate}
    \item membedakan kesalahan sintaks, kesalahan logika, dan program
          yang berhenti;
    \item memakai perintah DEBUG \texttt{R}, \texttt{T}, dan
          \texttt{D} serta \emph{single step} emu8086;
    \item menelusuri kesalahan perulangan dan kesalahan stack;
    \item memperbaiki kedua kesalahan tanpa menambah jenis kesalahan
          lain;
    \item memverifikasi perbaikan dengan penelusuran ulang.
\end{enumerate}
\end{tujuan}

\subsection{Dasar singkat}

\subsubsection*{Jenis kesalahan dan alat}

Kesalahan sintaks terdeteksi assembler dan membuat berkas
\texttt{.OBJ} tidak terbentuk. Kesalahan logika membuat program tetap
berjalan, tetapi hasilnya salah. Program yang berhenti tidak kembali
ke DOS, biasanya karena perulangan tak berujung atau stack tidak
seimbang. Perbedaannya dirangkum pada Tabel~\ref{tab:m11-jenis}.

\begin{table}[H]
    \centering
    \caption{Tiga jenis kesalahan dan alat utamanya.}
    \label{tab:m11-jenis}
    \small
    \begin{tabular}{@{}lll@{}}
        \toprule
        \textbf{Jenis} & \textbf{Gejala} & \textbf{Alat utama} \\
        \midrule
        Sintaks & pesan TASM, tidak ada \texttt{.OBJ} &
        daftar kesalahan assembler \\
        Logika & program selesai, hasil salah &
        \emph{single step}, register \\
        Berhenti & program tidak kembali ke DOS &
        \emph{single step}, pemeriksaan stack \\
        \bottomrule
    \end{tabular}
\end{table}

DEBUG menyediakan perintah \texttt{R} untuk menampilkan register dan
flag, \texttt{T} untuk menjalankan satu instruksi, dan \texttt{D}
untuk menampilkan isi memori. emu8086 menyediakan tombol
\emph{Single Step} serta jendela register, flag, dan memori yang
diperbarui setiap instruksi.

\subsubsection*{Kesalahan perulangan}

Program pertama dimaksudkan menjumlahkan 1 sampai 5. Di dalam badan
perulangan, \texttt{INC CX} menaikkan pencacah, sedangkan
\texttt{LOOP} selalu menurunkannya. Kedua perubahan itu saling
meniadakan, sehingga \texttt{CX} tidak pernah mencapai nol dan
perulangan berjalan terus. Jejak beberapa iterasi ditunjukkan pada
Tabel~\ref{tab:m11-loop}.

\begin{table}[H]
    \centering
    \caption{Jejak kesalahan perulangan pada program pertama.}
    \label{tab:m11-loop}
    \small
    \begin{tabular}{@{}clll@{}}
        \toprule
        \textbf{Iterasi} & \textbf{Instruksi} & \texttt{AX} &
        \texttt{CX} \\
        \midrule
        1 & \texttt{ADD AX,CX} & 5 & 5 \\
        1 & \texttt{INC CX} & 5 & 6 \\
        1 & \texttt{LOOP ULANG} & 5 & 5, kembali ke awal \\
        2 & \texttt{ADD AX,CX} & 10 & 5 \\
        2 & \texttt{INC CX} & 10 & 6 \\
        2 & \texttt{LOOP ULANG} & 10 & 5, tidak pernah nol \\
        \bottomrule
    \end{tabular}
\end{table}

Perbaikan dilakukan dengan menghapus \texttt{INC CX}. Setelah itu
\texttt{CX} menurun dari 5 sampai 0 dan hasilnya 15. Hasil itu perlu
dicetak sebagai dua digit, sehingga rutin pembagian pada modul
sebelumnya dapat dipakai.

\subsubsection*{Kesalahan stack}

Program kedua memanggil prosedur \texttt{HITUNG}. Prosedur mendorong
tiga register dengan \texttt{PUSH}, tetapi hanya memulihkan dua
dengan \texttt{POP}. Saat \texttt{RET} dijalankan, puncak stack
berisi nilai \texttt{AX} lama, bukan alamat kembali. Prosesor
melompat ke alamat yang salah dan program berhenti. Susunan stack
ditunjukkan pada Tabel~\ref{tab:m11-stack}.

\begin{table}[H]
    \centering
    \caption{Susunan stack pada program kedua.}
    \label{tab:m11-stack}
    \small
    \begin{tabular}{@{}lll@{}}
        \toprule
        \textbf{Tahap} & \textbf{Puncak stack} & \textbf{Keterangan} \\
        \midrule
        sesudah \texttt{CALL} & alamat kembali &
        disimpan oleh \texttt{CALL} \\
        sesudah tiga \texttt{PUSH} & \texttt{AX}, \texttt{BX},
        \texttt{CX} & tiga word tersimpan \\
        sesudah dua \texttt{POP} & \texttt{AX} lama &
        satu word belum dibuang \\
        \texttt{RET} & \texttt{AX} lama dianggap alamat &
        lompat ke alamat salah \\
        \bottomrule
    \end{tabular}
\end{table}

Perbaikan dilakukan dengan menambahkan \texttt{POP AX} sebelum
\texttt{RET} sehingga jumlah \texttt{PUSH} dan \texttt{POP}
seimbang. Verifikasi dilakukan dengan menelusuri ulang sampai program
kembali ke DOS dan pesan penutup tercetak.

\begin{rujukan}
Rujukan singkat: S'to, Bab XXVII; Gunawan, Bagian 4.3; buku ajar,
\texttt{chapters/16} serta berkas \texttt{bug-loop.asm},
\texttt{bug-loop-fix.asm}, \texttt{bug-stack.asm}, dan
\texttt{bug-stack-fix.asm}.
\end{rujukan}

\subsection{Alat dan bahan}

\begin{itemize}
    \item Komputer dengan emu8086 terpasang.
    \item DOSBox dengan TASM, TLINK, dan DEBUG.
    \item Berkas \texttt{code/m11-bug-loop.asm} dan
          \texttt{code/m11-bug-stack.asm}.
    \item Lembar jurnal praktikum.
\end{itemize}

\subsection{Contoh program}

Kedua program berikut sengaja mengandung satu kesalahan. Program
pertama tidak pernah selesai karena pencacah perulangannya dinaikkan.
Program kedua melompat ke alamat salah karena stack tidak seimbang.

\lstinputlisting[
    caption={Program ber-bug: perulangan tidak berhenti.},
    label={lst:m11-bug-loop}
]{code/m11-bug-loop.asm}

\lstinputlisting[
    caption={Program ber-bug: jumlah \texttt{PUSH} dan \texttt{POP}
    tidak seimbang.},
    label={lst:m11-bug-stack}
]{code/m11-bug-stack.asm}

\subsection{Langkah kerja}

\begin{enumerate}
    \item Siapkan folder kerja, lalu salin kedua berkas ber-bug ke
          dalamnya. Jangan mengubah berkas aslinya.
    \item Rakit dan tautkan program pertama. Jalankan di emu8086
          dengan \emph{Single Step}. Catat \texttt{AX} dan
          \texttt{CX} pada tiga iterasi pertama, lalu hentikan
          penelusuran. Di DEBUG, perintah \texttt{R} menampilkan
          register dan \texttt{T} menjalankan satu instruksi.
    \item Perbaiki program pertama dengan menghapus \texttt{INC CX}.
          Simpan sebagai \texttt{m11-bug-loop-\allowbreak fix.asm},
          rakit, tautkan, dan jalankan. Pastikan hasil 15 tercetak dan
          program kembali ke DOS.
    \item Rakit dan tautkan program kedua. Jalankan dan catat
          gejalanya. Telusuri dengan \emph{single step} sampai masuk
          ke prosedur \texttt{HITUNG}. Hitung \texttt{PUSH} dan
          \texttt{POP}, lalu periksa stack dengan \texttt{D} memakai
          \texttt{SS} dan \texttt{SP} yang dibaca dari \texttt{R}.
    \item Perbaiki program kedua dengan menambahkan \texttt{POP AX}
          sebelum \texttt{RET}. Simpan sebagai
          \texttt{m11-bug-stack-\allowbreak fix.asm}, rakit, tautkan,
          dan jalankan. Pastikan pesan \texttt{Selesai.} tercetak.
    \item Verifikasi kedua perbaikan dengan penelusuran ulang:
          periksa \texttt{AX} dan \texttt{CX} pada perbaikan pertama,
          serta \texttt{SP} sebelum \texttt{CALL} dan sesudah
          \texttt{RET} pada perbaikan kedua.
\end{enumerate}

\subsection{Tugas}

\begin{enumerate}
    \item Jelaskan perbedaan kesalahan sintaks, kesalahan logika, dan
          program yang berhenti. Berikan satu contoh untuk
          masing-masing.
    \item Jelaskan fungsi perintah DEBUG \texttt{R}, \texttt{T}, dan
          \texttt{D}.
    \item Jelaskan mengapa \texttt{INC CX} di dalam perulangan membuat
          program pertama tidak selesai. Sertakan jejak nilai
          \texttt{CX}.
    \item Jelaskan mengapa \texttt{RET} melompat ke alamat salah pada
          program kedua. Sertakan susunan stack sebelum \texttt{RET}.
    \item Jelaskan mengapa perbaikan harus diuji dengan penelusuran
          ulang, bukan hanya dengan membaca kode.
\end{enumerate}

\subsection{Hasil yang dikumpulkan}

\begin{enumerate}
    \item Kedua berkas ber-bug serta kedua berkas perbaikan:
          \begin{itemize}
              \item \texttt{m11-bug-loop-fix.asm}
              \item \texttt{m11-bug-stack-fix.asm}
          \end{itemize}
    \item Tangkapan layar penelusuran pada emu8086 atau DEBUG.
    \item Tabel jejak \texttt{AX}/\texttt{CX} dan tabel susunan stack
          sebelum serta sesudah perbaikan.
    \item Lembar jurnal berisi jawaban langkah kerja dan tugas.
\end{enumerate}
